\scanpage{227}
\margin{\figure{\begin{tikzpicture}\draw (0,.65)--(0,.12);\draw (0,.08) circle (.07);\end{tikzpicture}}
en vertu de la conjecture de Hodge}
D’autre part, \uline{quels} sont les foncteurs multiplicatifs de Motifs,
à valeurs dans les structures de Hodge,
\inserted{(i.e. \unsure{dans} les motifs sur $\mathbf C$ !)}
qui sont associés à des places archimédiennes ?
Plus généralement, dans quelle mesure un morphisme de préschémas
\inserted{$X\mathbin{—}Y$} est connu par le foncteur sur les motifs
\inserted{(loc. cit.)}, et dans quelle mesure un foncteur
\inserted{multiplicatif} peut-il être pris arbitrairement ???

C’est lié au Pb de \uline{construction des schémas de modules} pour les motifs
polarisés ; \unsure{si $M$} polarisé fidèle tel que
\inserted{($Y$ schéma sur \unsure{c. alg.})}
$Y\longrightarrow\mathrm{Modules}$ est une immersion
\inserted{(\unsure{resp. une}})
\deleted{\unsure{est}} alors \unsure{est|ça} \unsure{analogue}.
Kif kif si $Y$ est \uline{l’un} des préschémas ayant cette belle propriété.
Marchera en particulier si $Y$ est le spectre d’un corps.....
\margin{226}
